% =============================================================================
% REPRODUCIBILITY-AUDIT: A Single-Stack Survival Protocol for the GRPO Family
% =============================================================================
% Pillar 3 of the ZVF Program. Companion to the MIN-REPORT-RL position paper
% and the GRPO-Registry catalog.
%
% Self-contained draft. Compiles with plain article class; citations are TODO
% keys; bibliography is manual thebibliography stub.
% =============================================================================

\documentclass[11pt]{article}

\usepackage[utf8]{inputenc}
\usepackage[T1]{fontenc}
\usepackage[margin=1in]{geometry}
\usepackage{hyperref}
\usepackage{url}
\usepackage{booktabs}
\usepackage{amsmath}
\usepackage{amssymb}
\usepackage{xcolor}
\usepackage{enumitem}
\usepackage{microtype}
\usepackage{graphicx}

\renewcommand{\cite}[1]{\textcolor{red}{[\textsc{cite:}\,#1]}}
\newcommand{\todo}[1]{\textcolor{red}{\textbf{[TODO:\ #1]}}}
\newcommand{\zvf}{\textsc{Zvf}}
\newcommand{\gu}{\textsc{Gu}}
\newcommand{\minreport}{\textsc{Min-Report-RL}}
\newcommand{\grpo}{\textsc{GRPO}}
\newcommand{\audit}{\textsc{GRPO-Survival-Audit}}
\newcommand{\grpoReg}{\textsc{GRPO-Registry}}

\title{\audit{}: A Single-Stack Survival Protocol for the \grpo{} Family\\
\large Holding the Stack Fixed So Algorithmic Gains Can Be Seen}

\author{Arvind C R \and the ZVF Program\thanks{Pillar~3 of the ZVF Program.
Author list and affiliations \todo{finalize}.}}

\date{\today}

\begin{document}
\maketitle

% =============================================================================
\begin{abstract}
% =============================================================================
The \grpo{}-family literature reports head-to-head comparisons in which each
variant runs in its own trainer, sampler, and evaluation harness. Because the
``stack''---loss form, reference/KL handling, sampler precision, group-size
schedule, \zvf/\gu{} trajectory, held-out split, and parser---is a documented
flip lever \cite{minreportrl2026}, such comparisons confound algorithmic
innovation with implementation differences. This paper proposes \audit{}, a
single-stack reproducibility-audit protocol in which DAPO, GSPO, Dr.\grpo{},
MAD-\grpo{}, and other variance-mitigation variants are re-implemented as
minimal configuration overrides on one shared trainer. We specify the shared
stack, the pre-registration rules, the survival-verdict logic, the statistical
protocol, and the role of the companion \texttt{grpo-stackdiff} tool in
verifying that arms are stack-matched. We also report a small open-source pilot
on Qwen2.5-0.5B-Instruct (two seeds, T4 GPU) that demonstrates the protocol's
feasibility and produces concrete \zvf{} and held-out-delta observations. The
full-scale audit results are reported as placeholder tables to be filled once
compute is budgeted; the contribution is the protocol and the pilot, not a
final leaderboard.
\end{abstract}

% =============================================================================
\section{Introduction: The Need for a Survival Audit}
% =============================================================================

A typical GRPO-family paper compares a new variant against ``GRPO'' by running
each in a different codebase. The new variant may change the clip rule, the
advantage normalization, or the group-size schedule; the baseline may use a
different token mask, KL placement, sampler precision, or reward parser
\cite{minreportrl2026, zvfaudit2026}. The result is a comparison of two stacks
$S_A$ and $S_B$ that happen to be tagged ``A'' and ``B''. The algorithmic delta
and the stack delta are inseparable.

The \minreport{} position paper \cite{minreportrl2026} proposes a
seven-item minimum-reportable-stack standard so that future papers expose the
stack. The \grpoReg{} catalog \cite{grpoRegistry2026} records each variant's
self-declared delta relative to a baseline. This paper closes the loop: it
specifies a controlled, single-stack re-implementation protocol that asks,
for each variant, \emph{what fraction of its claimed gain survives when the
stack is held fixed?} We call the answer the variant's \emph{survival}.

\paragraph{What survival is not.} A low survival verdict does not mean a
variant is ``fake.'' It means the variant's published gain was partly or wholly
a stack effect under the audited stack. A high survival verdict is evidence
that the algorithmic idea is robust. Both outcomes are useful.

\paragraph{A live specimen.} While assembling this audit's own pilot, our
trainer shipped exactly the failure the protocol targets: a
\texttt{--loss drgrpo} flag whose commit message claimed it toggled the
advantage normalization, but which was never wired into the loss --- so a
six-arm GRPO-vs-Dr.GRPO panel silently compared GRPO with GRPO, and its
committed ``finding'' stood until the runner was read line by line. No
output-level check caught it: rewards, lengths, and ZVF traces all looked
plausible. The panel was invalidated, the artifacts preserved under
explicit \texttt{.invalid\_actually\_grpo.json} names, and the corrected
rerun (2026-07-11) reversed one of the two claims. Stack identity cannot
be certified from results; it must be certified from the stack
(\S\ref{sec:stackdiff}).

\paragraph{Contributions.}
\begin{enumerate}[leftmargin=1.6em, itemsep=2pt]
  \item A concrete single-stack audit protocol with shared-stack constraints,
    arm definitions, and pre-registration rules (\S\ref{sec:design}).
  \item A survival-verdict logic and statistical protocol that treat seed as
    the independent unit (\S\ref{sec:verdict}--\ref{sec:stats}).
  \item Integration with \texttt{grpo-stackdiff} to verify stack identity
    across arms before any result is trusted (\S\ref{sec:stackdiff}).
  \item A small open-source pilot (\S\ref{sec:pilot}) that exercises the
    protocol on four arms and shows that the required telemetry is obtainable
    from a standard GRPO trainer.
\end{enumerate}

% =============================================================================
\section{Audit Design}
\label{sec:design}
% =============================================================================

\paragraph{Shared stack.} Every arm uses the same:
\begin{itemize}[leftmargin=1.4em, itemsep=2pt]
  \item base checkpoint and tokenizer revision,
  \item LoRA target modules, rank, alpha, and dropout (if LoRA is used),
  \item optimizer, learning-rate schedule, and gradient clipping,
  \item rollout engine (vLLM), decoding parameters, sampler/trainer precision,
  \item reference-policy and KL handling (placement, coefficient, schedule,
    estimator),
  \item token mask and advantage normalization (unless the variant explicitly
    changes one of these),
  \item reward parser and held-out evaluation harness,
  \item prompt distribution and train/test split.
\end{itemize}
The only permitted differences are the variant's declared algorithmic deltas.

\paragraph{Arms.} The primary arms are:
\begin{itemize}[leftmargin=1.4em, itemsep=2pt]
  \item \textbf{GRPO (baseline):} vanilla per-token GRPO with the shared stack.
  \item \textbf{DAPO:} asymmetric clip-higher + dynamic sampling
    \cite{dapo2025}.
  \item \textbf{GSPO:} sequence-level importance-sampling ratio
    \cite{gspo2025}.
  \item \textbf{Dr.\grpo{}:} length-normalized advantage/reward terms
    \cite{drgrpo2025}.
  \item \textbf{MAD-\grpo{}:} multi-agent / diversity bonus
    \cite{madgrpo2025}.
\end{itemize}
Optional secondary arms (AERO \cite{aero2024}, CPPO \cite{cppo2024}, NGRPO
\cite{ngrpo2025}, Scaf-\grpo{} \cite{scafgrpo2025}) are added if compute
permits; each is implemented as a single hook on the same shared trainer.

\paragraph{Pre-registration.} Before any held-out score is observed, the audit
fixes: the arm list, the shared stack specification, the seed sweep, the
held-out split, the checkpoint-selection rule, the survival thresholds, and the
primary analysis script. Pre-registration prevents the selection-by-reward
bias that the audit is designed to diagnose \cite{zvfaudit2026}.

\paragraph{Hook interface.} A variant is added by implementing a small hook
object with the following interface:
\begin{verbatim}
class VariantHook:
    def on_group_rewards(self, rewards: Tensor) -> Tensor:
        # optionally transform per-prompt group rewards
        return rewards
    def on_advantage(self, advantage: Tensor) -> Tensor:
        # optionally transform advantages
        return advantage
    def on_loss(self, loss_terms: Tensor, mask: Tensor) -> Tensor:
        # optionally modify loss terms
        return loss_terms
    def group_size_rule(self, step: int, zvf: float) -> int:
        # return G for this step; default is fixed G
        return self.G
\end{verbatim}
All hooks are loaded into the same trainer binary, so infrastructure differences
are impossible.

% =============================================================================
\section{Survival Verdicts}
\label{sec:verdict}
% =============================================================================

For each variant we report two numbers:
\begin{itemize}[leftmargin=1.4em, itemsep=2pt]
  \item \textbf{Published $\Delta$:} the gain claimed in the original paper on
    its own stack.
  \item \textbf{Controlled $\Delta$:} the gain in the shared stack vs.\ the
    shared GRPO baseline, with a 95\% confidence interval across seeds.
\end{itemize}
The survival verdict is then:
\begin{center}
\small
\begin{tabular}{@{}ll@{}}
\toprule
\textbf{Verdict} & \textbf{Rule}\\
\midrule
Survives & Controlled CI lower bound $>50\%$ of published $\Delta$.\\
Shrinks & Controlled CI overlaps $(0, \text{published }\Delta)$ but does not
  fully cover it.\\
Disappears & Controlled CI includes zero.\\
Reverses & Controlled CI is entirely negative.\\
\bottomrule
\end{tabular}
\end{center}
The $50\%$ survival threshold is a pre-registered convention; sensitivity
analyses at $25\%$ and $75\%$ are reported in an appendix.

% =============================================================================
\section{Statistical Protocol}
\label{sec:stats}
% =============================================================================

The independent unit is a seed, not a training step. Training rewards are
autocorrelated within a trace, so inferential claims are made only across the
seed sweep. For each metric we report the mean and standard error across seeds,
with bootstrap 95\% confidence intervals. We control the false-discovery rate
across the pre-registered survival tests using Benjamini--Hochberg. Single-seed
arms are reported as descriptive only.

The minimum detectable effect at the chosen seed count $S$ is reported for the
primary held-out metric. If the published $\Delta$ is smaller than the MDE, the
audit cannot claim disappearance; it can only claim \texttt{INCONCLUSIVE}.

% =============================================================================
\section{Using \texttt{grpo-stackdiff} to Verify Stack Identity}
\label{sec:stackdiff}
% =============================================================================

Before interpreting any survival result, we verify that the arms are
stack-matched. The companion \texttt{grpo-stackdiff} CLI
\cite{minreportrl2026} takes the \minreport{} manifest of each arm and reports
whether any unapproved lever differs. The audit proceeds only if every pairwise
comparison returns \texttt{STACK\_MATCHED} (or \texttt{STACK\_MATERIAL} for
levers the protocol explicitly tolerates).

If a hook inadvertently changes a non-target lever---for example, a dynamic
sampling rule that also changes the effective tokenizer chat template because
it filters on string prefixes---\texttt{grpo-stackdiff} flags it as a
DISTRIBUTIONAL or SEMANTIC\_OBJECTIVE diff. The hook must then be rewritten so
that only the declared delta remains. This check is the difference between a
``single-trainer'' audit and a genuinely ``single-stack'' audit.

% =============================================================================
\section{Primary Results Tables}
% =============================================================================

Table~\ref{tab:survival} is the headline table. All numeric cells are
placeholders to be filled once the full-scale audit corpus is collected.
Table~\ref{tab:telemetry} reports the \minreport{} item-4 telemetry for each
arm.

\begin{table}[t]
\centering
\caption{\textbf{Controlled single-stack survival audit.} Each variant is a
  minimal hook on the shared trainer. ``Published $\Delta$'' is the original
  paper's claim; ``Controlled $\Delta$'' is the shared-stack gain vs.\ the
  shared GRPO baseline. Cells are \todo{fill from full-scale audit corpus}.}
\label{tab:survival}
\small
\resizebox{\textwidth}{!}{%
\begin{tabular}{@{}lcccccc@{}}
\toprule
\textbf{Variant} & \textbf{Defining delta} & \textbf{Published $\Delta$} &
\textbf{Controlled last-10} & \textbf{Controlled held-out} &
\textbf{Controlled $\Delta$ (95\% CI)} & \textbf{Survival}\\
\midrule
GRPO (baseline) & --- & --- & \todo{} & \todo{} & $0$ (ref.) & --- \\
DAPO     & clip-higher + dyn.\ sampling & \todo{} & \todo{} & \todo{} & \todo{} & \todo{} \\
GSPO     & sequence-level IS ratio     & \todo{} & \todo{} & \todo{} & \todo{} & \todo{} \\
Dr.\grpo{} & length/normalization fix   & \todo{} & \todo{} & \todo{} & \todo{} & \todo{} \\
MAD-\grpo{} & multi-agent / diversity   & \todo{} & \todo{} & \todo{} & \todo{} & \todo{} \\
\bottomrule
\end{tabular}}
\end{table}

\begin{table}[t]
\centering
\caption{\textbf{Per-arm telemetry.} Mean \zvf{} and \gu{} late in training,
  collapse behavior, and group-size schedule under the shared stack. Cells are
  \todo{fill from full-scale audit corpus}.}
\label{tab:telemetry}
\small
\resizebox{\textwidth}{!}{%
\begin{tabular}{@{}lccccc@{}}
\toprule
\textbf{Variant} & \textbf{Mean \zvf{} @ step 25} & \textbf{Mean \gu{} @ step 25} &
\textbf{Collapse rate} & \textbf{Median steps to collapse} &
\textbf{Group-size schedule}\\
\midrule
GRPO (baseline) & \todo{} & \todo{} & \todo{} & \todo{} & fixed $G$ \\
DAPO     & \todo{} & \todo{} & \todo{} & \todo{} & dynamic \\
GSPO     & \todo{} & \todo{} & \todo{} & \todo{} & fixed \todo{} \\
Dr.\grpo{} & \todo{} & \todo{} & \todo{} & \todo{} & fixed \todo{} \\
MAD-\grpo{} & \todo{} & \todo{} & \todo{} & \todo{} & \todo{} \\
\bottomrule
\end{tabular}}
\end{table}

% =============================================================================
\section{Open-Source Pilot}
\label{sec:pilot}
% =============================================================================

We ran a small pilot on a single T4 GPU to validate that the required telemetry
can be extracted from a standard GRPO trainer and that the protocol produces
informative stack-controlled observations. The pilot is \emph{not} the
full-scale audit; it uses a small model and only two seeds.

\paragraph{Setup.} Model: Qwen/Qwen2.5-0.5B-Instruct. Two seeds. Arms: GRPO,
Dr.\grpo{}, DAPO, and an adaptive-$G$ GRPO baseline. The shared stack is a
single open-source trainer with the same tokenizer, sampler, reward parser, and
held-out split across arms.

\paragraph{Pilot results.} Table~\ref{tab:pilot} shows mean held-out delta,
mean \zvf{}, and mean rollouts per arm. Even at this scale, the stack-controlled
observations are informative: DAPO reports zero \zvf{} but spends more rollouts
(174 vs.\ 120), while adaptive-$G$ matches Dr.\grpo{}'s held-out delta with a
different \zvf{} and rollout budget.

\begin{table}[t]
\centering
\caption{\textbf{Pilot results (Qwen2.5-0.5B-Instruct, 2 seeds, T4).} Mean
  held-out $\Delta$, mean \zvf{}, and mean rollouts per arm.}
\label{tab:pilot}
\small
\begin{tabular}{@{}lccc@{}}
\toprule
\textbf{Arm} & \textbf{Held-out $\Delta$} & \textbf{Mean \zvf{}} &
\textbf{Mean rollouts}\\
\midrule
GRPO (baseline)          & $0.500$ & $0.250$ & $120$ \\
Dr.\grpo{}                & $0.575$ & $0.267$ & $120$ \\
DAPO                     & $0.550$ & $0.000$ & $174$ \\
GRPO adaptive-$G$        & $0.575$ & $0.233$ & $186$ \\
\bottomrule
\end{tabular}
\end{table}

\paragraph{Interpretation.} The pilot shows that the same held-out delta can
arrive through very different signal-budget paths. DAPO eliminates zero-variance
groups at the cost of more rollouts; adaptive-$G$ achieves the same delta as
Dr.\grpo{} with a higher rollout budget. A head-to-head that reports only final
reward would miss these trade-offs; the \audit{} protocol surfaces them
explicitly.

% =============================================================================
\section{Limitations and Extensions}
% =============================================================================

\paragraph{Scope.} The audit tests survival under one shared stack. A variant
that survives here may still fail under a different stack; the audit does not
claim universal robustness. It claims only that the published gain is not
purely an artifact of the original stack.

\paragraph{Faithfulness of re-implementation.} A variant may depend on an
unreported implementation detail. If a careful re-implementation cannot
reproduce the original result from the paper, that is itself a finding about
reporting quality, but it is not a finding about the method's intrinsic value.
We report both readings and publish every arm's full \minreport{} block.

\paragraph{Extensions.} The protocol generalizes to any family of methods that
are small deltas on a base loop. The same hook interface can accommodate
future variants; the same \texttt{grpo-stackdiff} check can verify that each
new hook changes only its declared levers.

% =============================================================================
\section{Conclusion}
% =============================================================================

The \grpo{} family needs a way to separate algorithmic gains from stack
effects. The \minreport{} standard tells authors what to report; the
\grpoReg{} catalog records each variant's declared delta; and the \audit{}
protocol reported here holds the stack fixed so that those deltas can be tested
fairly. The immediate next step is to run the full-scale, multi-seed audit on a
capable model and fill the placeholder tables. Until then, the protocol, the
pilot, and the \texttt{grpo-stackdiff} verification step are the contribution:
they make it possible, for the first time, to ask of a GRPO-family result not
``did it win?'' but ``did it survive?''

% =============================================================================
\begin{thebibliography}{9}

\bibitem{minreportrl2026}
\todo{\minreport{} position paper. Cite arXiv/venue version once assigned.}

\bibitem{grpoRegistry2026}
\todo{\grpoReg{} catalog paper. Cite arXiv/venue version once assigned.}

\bibitem{zvfaudit2026}
\todo{ZVF Program v1 audit paper. Cite arXiv/venue version once assigned.}

\bibitem{dapo2025}
\todo{DAPO paper. Fill authors/venue/year.}

\bibitem{gspo2025}
\todo{GSPO paper. Fill authors/venue/year.}

\bibitem{drgrpo2025}
\todo{Dr.\grpo{} paper. Fill authors/venue/year.}

\bibitem{madgrpo2025}
\todo{MAD-\grpo{} paper. Fill authors/venue/year.}

\bibitem{aero2024}
\todo{AERO paper. Fill authors/venue/year.}

\bibitem{cppo2024}
\todo{CPPO paper. Fill authors/venue/year.}

\bibitem{ngrpo2025}
\todo{NGRPO paper. Fill authors/venue/year.}

\bibitem{scafgrpo2025}
\todo{Scaf-\grpo{} paper. Fill authors/venue/year.}

\end{thebibliography}

\end{document}
